\documentclass[tikz,border=4pt]{standalone}
\usepackage{ctex}
\usepackage{amsmath}
\usepackage{pgfplots}
\pgfplotsset{compat=1.18}
\usetikzlibrary{calc, arrows.meta}

\begin{document}
\begin{tikzpicture}[scale=1.0, thick]
  % 坐标轴
  \draw[->] (-0.5,0) -- (5.0,0) node[right] {$x$};
  \draw[->] (0,-0.5) -- (0,5.0) node[above] {$y$};
  \node[below left] at (0,0) {$O$};

  % 刻度
  \foreach \x in {1,2,3,4} {
    \draw (\x, 2pt) -- (\x, -2pt) node[below] {\small$\x$};
  }
  \foreach \y in {1,2,3,4} {
    \draw (2pt, \y) -- (-2pt, \y) node[left] {\small$\y$};
  }

  % 直线1: 2x + y = 7, i.e. y = 7 - 2x (过 (1,5) 和 (3.5,0))
  \draw[blue, thick] (1.0, 5.0) -- (3.5, 0.0) node[right, blue] {$2x+y=7$};

  % 直线2: x - y = 2, i.e. y = x - 2 (过 (2,0) 和 (4,2))
  \draw[red, thick] (2.0, 0.0) -- (4.5, 2.5) node[right, red] {$x-y=2$};

  % 交点 (3,1)
  \fill (3,1) circle (2.5pt);
  \node[above right] at (3,1) {$(3,\,1)$};

  % 辅助虚线
  \draw[gray, dashed] (3,0) -- (3,1);
  \draw[gray, dashed] (0,1) -- (3,1);
  \node[below] at (3,0) {\small$3$};
  \node[left] at (0,1) {\small$1$};

  % 说明
  \node[below] at (2.5,-0.9) {\small 方程组的解 $(3,1)$ = 两直线的交点};
\end{tikzpicture}
\end{document}
